\section{Заключение}

В работе была выполнена концептуальная репликация подхода Sims и Zha на российских квартальных данных. Построенная VAR/SVAR-модель позволяет разделить шоки денежного предложения, спроса на деньги и прочие макроэкономические возмущения.

Функции импульсного отклика в целом согласуются с базовой логикой денежно-кредитной трансмиссии: шок денежного предложения отражается в динамике ставки, денежной массы, выпуска и цен. При этом результаты не дают оснований напрямую связывать макроэкономические спады только с неожиданными шоками денежно-кредитной политики. Это соответствует выводу Sims и Zha: большая часть изменений в денежных переменных отражает систематическую реакцию политики, а роль неожиданных денежных шоков в циклических колебаниях ограничена.

Контрфактический анализ с фиксированными $M2$ и $R$ показывает, что систематическая часть политики влияет на реакцию экономики на шоки, но эти результаты зависят от идентифицирующих ограничений.

Основной вывод состоит в том, что денежно-кредитная политика может быть важным каналом передачи шоков, но по результатам данной модели её нельзя считать единственным или главным источником экономических спадов.
